\chapter{Bounded concurrency and connection lifecycle}
\chapterintro{Correct parsing is not enough if one silent client can consume every worker. The final server makes capacity, waiting, persistence, overload, and shutdown explicit.}

\begin{objectives}
\begin{itemize}
  \item Explain the accept loop, finite connection queue, and fixed workers.
  \item Handle per-connection parsing, timeouts, keep-alive, and maximum requests.
  \item Define overload behavior instead of permitting unbounded memory growth.
  \item Stop accepting and release resources during graceful shutdown.
\end{itemize}
\end{objectives}

\section{From serial to concurrent}

The early snapshots handle one accepted connection at a time. If a client pauses,
new clients wait even when CPU and database capacity are available. Stage 7 adds
a fixed set of worker threads and a finite queue:

\codefile{../code/stages/07_bounded_concurrency/server.py}

The accept loop does little work: accept, enqueue, or reject. Each worker takes a
connection, applies a timeout, processes it, closes it, and marks the queue item
complete. Four workers permit four blocking conversations; a queue of sixteen
allows a short burst without allocating unlimited threads.

\section{Bounded is a correctness property}

An unbounded queue converts overload into growing memory use and increasing
latency. An unbounded thread-per-connection design can exhaust memory, scheduler
capacity, descriptors, and database connections. A bounded service reaches a
known saturation point and rejects excess work predictably.

The final defaults are sixteen workers and a 64-connection application queue.
When the queue is full, the accept loop sends a small 503 problem and closes the
connection. The response is best effort: a peer may disappear before reading it.

Capacity values are not universal recommendations. Production sizing requires
measurement of workload, latency, memory, downstream capacity, and failure
behavior. Here they are visible teaching defaults.

\section{One parser per connection}

Each worker constructs a \code{RequestParser} when it begins a connection. The
parser buffer survives across receive calls and, for keep-alive, across requests.
It is never shared between connections. This avoids locks inside the parser and
prevents bytes from different clients mixing.

The worker loop asks the parser for a request before receiving more data. This
detail matters when the previous receive already contained two requests. It
processes buffered work without waiting for network input that may never come.

\section{Persistent but finite}

HTTP/1.1 persistence is enabled by default. The server closes when:

\begin{itemize}
  \item the client sends \code{Connection: close};
  \item keep-alive is disabled in configuration;
  \item 100 requests have used the connection;
  \item parsing fails, input times out, the peer closes, or socket I/O fails;
  \item server shutdown begins.
\end{itemize}

Each response states \code{Connection: keep-alive} or \code{close}. The five-
second socket timeout bounds both a partial request and idle time before the next
request. A production server would often use separate header, body, idle, write,
and total deadlines.

\section{Read the final server}

\codefile{../code/final/tinyhttp/server.py}

\subsection{Readiness and ephemeral ports}

When configured with port zero, the operating system selects a free port. The
server records its bound address and sets a readiness event after listening and
starting workers. Integration tests can start a server thread, wait deterministically,
and connect to the chosen port without guessing.

\subsection{Timeout and protocol errors}

An \code{HTTPError} becomes a problem response and closes the connection so no
ambiguous bytes survive. A timeout with buffered data can receive 408; an idle
connection simply closes. Socket errors are normal peer-failure outcomes and do
not crash the accept loop.

\subsection{Shutdown}

\code{shutdown} sets a stop event, closes the listener to wake accept, places
sentinels for workers, and joins them within a bound. Workers also observe the
event in their loops. This is modest graceful behavior: it stops new acceptance
and lets current operations finish or time out. It does not provide a configurable
long drain period or connection migration.

\section{Threads and SQLite}

Python threads are appropriate for this teaching server because socket and
SQLite waits release execution opportunities and blocking control flow remains
readable. They do not make CPU-bound Python code scale linearly across cores.

The Task repository opens a short-lived SQLite connection for each operation.
Connections are not shared across workers. SQLite serializes writes as needed,
and a busy timeout gives brief contention time to resolve. WAL mode improves
reader/writer coexistence, but SQLite remains one embedded database with finite
write capacity.

\section{Alternatives}

An event loop can manage many sockets with fewer threads by running non-blocking
operations and resuming tasks when ready. ASGI frameworks use asynchronous
interfaces for such workloads. Process pools or multiple server processes can
use more CPU cores and isolate failures. Mature servers combine strategies.

The important lesson is not ``threads are best.'' It is that every concurrency
model has ownership, scheduling, cancellation, backpressure, and shutdown rules.
Our bounded thread model makes those rules small enough to inspect.

\begin{securityboundary}
Timeouts and queues reduce resource exhaustion but do not make this server safe
for hostile Internet traffic. Production parsers, TLS stacks, operating-system
tuning, rate limiting, monitoring, patch processes, and defense in depth remain
necessary.
\end{securityboundary}

\summarytitle

The server accepts into a finite queue and uses a fixed worker set. Each
connection owns its parser and is constrained by timeout and request count.
Overload receives 503, protocol errors close ambiguous streams, and shutdown
releases listener and workers. Bounds turn failure from accidental exhaustion
into explicit behavior.

\reviewtitle
\begin{enumerate}
  \item Why is a fixed worker count insufficient without a bounded queue?
  \item Why must a parser buffer be connection-local?
  \item Which conditions end keep-alive in this implementation?
  \item What work remains for a production-grade graceful shutdown?
\end{enumerate}

\begin{exercise}[Capacity timeline]
With two workers and a queue of three, draw the state after six clients connect
and each worker blocks on a partial header. Identify which clients are active,
queued, or rejected. Repeat after one worker times out.
\end{exercise}

\begin{exercise}[Concurrency test]
Write an integration test that starts the server on port zero and sends 50
concurrent GET requests. Collect every status with a deadline. The assertion
should allow success and intentional overload, but no hangs or malformed
responses.
\end{exercise}

